\documentclass[12pt,a4paper]{article}
\usepackage[utf8]{inputenc}
\usepackage[italian]{babel}
\usepackage{geometry}
\usepackage{graphicx}
\usepackage{listings}
\usepackage{xcolor}
\usepackage{hyperref}
\usepackage{booktabs}
\usepackage{array}
\usepackage{amsmath}
\usepackage{longtable}
\usepackage{multirow}
\usepackage{tikz}
\usetikzlibrary{arrows.meta, positioning, shapes.geometric}

\geometry{margin=2.5cm}

\definecolor{codegreen}{rgb}{0,0.6,0}
\definecolor{codegray}{rgb}{0.5,0.5,0.5}
\definecolor{codepurple}{rgb}{0.58,0,0.82}
\definecolor{backcolour}{rgb}{0.95,0.95,0.92}
\definecolor{okgreen}{rgb}{0.0,0.5,0.0}
\definecolor{warnOrange}{rgb}{0.9,0.6,0.0}
\definecolor{errorRed}{rgb}{0.8,0.0,0.0}

\lstdefinestyle{javastyle}{
    backgroundcolor=\color{backcolour},
    commentstyle=\color{codegreen},
    keywordstyle=\color{blue},
    numberstyle=\tiny\color{codegray},
    stringstyle=\color{codepurple},
    basicstyle=\ttfamily\footnotesize,
    breakatwhitespace=false,
    breaklines=true,
    captionpos=b,
    keepspaces=true,
    numbers=left,
    numbersep=5pt,
    showspaces=false,
    showstringspaces=false,
    showtabs=false,
    tabsize=2,
    language=Java
}

\newcommand{\ok}{\textcolor{okgreen}{\checkmark}}
\newcommand{\warn}{\textcolor{warnOrange}{$\sim$}}
\newcommand{\bad}{\textcolor{errorRed}{$\times$}}

\title{\textbf{Analisi GRASP Completa}\\[0.3cm]
\large Use Case ``Crea Ordine'' - Habibi Yalla\\[0.2cm]
\normalsize Valutazione delle Classi Coinvolte}
\author{Progetto ISPW 2026}
\date{Febbraio 2026}

\begin{document}

\maketitle

\tableofcontents
\newpage

%==============================================================================
\section{Panoramica dell'Architettura}

Il caso d'uso ``Crea Ordine'' segue l'architettura \textbf{BCE} (Boundary-Control-Entity) con i seguenti componenti:

\begin{figure}[h]
\centering
\begin{tikzpicture}[
    node distance=1.5cm,
    box/.style={rectangle, draw, rounded corners, minimum width=3cm, minimum height=1cm, align=center},
    arrow/.style={-{Stealth}, thick}
]
    % Boundary
    \node[box, fill=blue!20] (gui) {CreaOrdine\\GUIController};
    \node[box, fill=blue!20, right=of gui] (cli) {CreaOrdine\\CLIController};
    
    % Facade
    \node[box, fill=yellow!30, below=1.5cm of $(gui)!0.5!(cli)$] (facade) {CreaOrdine\\Facade};
    
    % Control
    \node[box, fill=green!20, below left=1.5cm and 0.5cm of facade] (ctrl) {CreaOrdine\\Controller};
    \node[box, fill=green!20, below right=1.5cm and 0.5cm of facade] (voucher) {UsaVoucher\\Controller};
    
    % Entity
    \node[box, fill=red!20, below=1.5cm of ctrl] (ordine) {Ordine};
    \node[box, fill=red!20, right=of ordine] (food) {Food};
    \node[box, fill=red!20, right=of food] (voucherE) {Voucher};
    
    % Pure Fabrication
    \node[box, fill=purple!20, right=3cm of ctrl] (mapper) {Mappers};
    \node[box, fill=purple!20, right=of mapper] (factory) {FoodFactory};
    
    % Arrows
    \draw[arrow] (gui) -- (facade);
    \draw[arrow] (cli) -- (facade);
    \draw[arrow] (facade) -- (ctrl);
    \draw[arrow] (ctrl) -- (voucher);
    \draw[arrow] (ctrl) -- (ordine);
    \draw[arrow] (ctrl) -- (food);
    \draw[arrow] (voucher) -- (voucherE);
    \draw[arrow, dashed] (ctrl) -- (mapper);
    \draw[arrow, dashed] (ctrl) -- (factory);
    
    % Labels
    \node[above=0.3cm of gui, xshift=1.5cm] {\textbf{BOUNDARY}};
    \node[left=0.3cm of facade] {\textbf{FACADE}};
    \node[left=0.3cm of ctrl] {\textbf{CONTROL}};
    \node[left=0.3cm of ordine] {\textbf{ENTITY}};
    \node[right=0.3cm of factory] {\textbf{PURE FABRICATION}};
    
\end{tikzpicture}
\caption{Architettura BCE del caso d'uso ``Crea Ordine''}
\end{figure}

\subsection{Classi Coinvolte}

\begin{table}[h]
\centering
\begin{tabular}{|l|c|c|l|}
\hline
\textbf{Classe} & \textbf{Layer} & \textbf{LOC} & \textbf{Responsabilità} \\
\hline
CreaOrdineGUIController & Boundary & 485 & Interazione utente JavaFX \\
CreaOrdineCLIController & Boundary & 375 & Interazione utente CLI \\
CreaOrdineFacade & Facade & 72 & Punto di accesso semplificato \\
CreaOrdineController & Control & 235 & Business logic ordine \\
UsaVoucherController & Control & 142 & Business logic voucher \\
Ordine & Entity & 251 & Modello dati ordine \\
Food (+ sottoclassi) & Entity & $\sim$300 & Modello dati prodotti \\
Voucher (+ sottoclassi) & Entity & $\sim$200 & Modello dati voucher \\
FoodMapper & Pure Fab. & 66 & Conversione Food$\to$Bean \\
VoucherMapper & Pure Fab. & 55 & Conversione Voucher$\to$Bean \\
OrdineMapper & Pure Fab. & 46 & Conversione Ordine$\to$Bean \\
FoodFactory & Pure Fab. & 75 & Creazione Food + Decorator \\
\hline
\end{tabular}
\caption{Riepilogo classi analizzate}
\end{table}

%==============================================================================
\section{Analisi per Principio GRASP}

%------------------------------------------------------------------------------
\subsection{Information Expert}

\begin{quote}
\textit{``Assegna la responsabilità alla classe che possiede le informazioni necessarie per adempierla.''}
\end{quote}

\subsubsection{Valutazione per Classe}

\begin{table}[h]
\centering
\begin{tabular}{|l|c|p{8cm}|}
\hline
\textbf{Classe} & \textbf{Voto} & \textbf{Motivazione} \\
\hline
Ordine & \ok & Calcola \texttt{getTotale()}, \texttt{getSubtotale()}, \texttt{getDurataTotale()} usando i propri dati \\
\hline
Food & \ok & Gestisce \texttt{getCosto()}, \texttt{getDurata()} internamente \\
\hline
Voucher & \ok & \texttt{calcolaSconto()} implementato nelle sottoclassi che conoscono le regole \\
\hline
CreaOrdineController & \ok & Coordina ma delega calcoli all'Entity \\
\hline
CreaOrdineGUIController & \ok & Gestisce solo UI, non calcoli di business \\
\hline
\end{tabular}
\caption{Information Expert - Valutazione}
\end{table}

\textbf{Esempio positivo} - Il calcolo del totale è nell'Entity \texttt{Ordine}:

\begin{lstlisting}[style=javastyle, caption={Information Expert in Ordine}]
public double getTotale() {
    // L'Ordine conosce i propri prodotti e voucher
    return getSubtotale() - getSconto();
}

public double getSubtotale() {
    double totale = 0;
    for (Food f : prodotti) {
        totale += f.getCosto(); // Delega a Food
    }
    return totale;
}
\end{lstlisting}

%------------------------------------------------------------------------------
\subsection{Creator}

\begin{quote}
\textit{``Assegna la responsabilità di creare un oggetto A alla classe B se B contiene/aggrega A, registra A, usa strettamente A, o ha i dati per inizializzare A.''}
\end{quote}

\subsubsection{Valutazione per Classe}

\begin{table}[h]
\centering
\begin{tabular}{|l|c|p{8cm}|}
\hline
\textbf{Classe} & \textbf{Voto} & \textbf{Motivazione} \\
\hline
FoodFactory & \ok & Centralizza creazione Food e Decorator (Pure Fabrication giustifica) \\
\hline
OrdineLazyFactory & \ok & Crea Ordine gestendo cache e persistenza \\
\hline
CreaOrdineController & \warn & Delega a Factory invece di creare direttamente (accettabile) \\
\hline
CreaOrdineFacade & \ok & Crea CreaOrdineController (lo contiene) \\
\hline
\end{tabular}
\caption{Creator - Valutazione}
\end{table}

\textbf{Esempio} - CreaOrdineFacade crea CreaOrdineController:

\begin{lstlisting}[style=javastyle, caption={Creator in CreaOrdineFacade}]
public CreaOrdineFacade(String tokenKey) throws MissingAuthorizationException {
    this.sessionUser = SessionManager.getInstance().getSessionUserByTokenKey(tokenKey);
    // Il Facade CONTIENE il Controller, quindi lo crea (GRASP Creator)
    this.controller = new CreaOrdineController();
}
\end{lstlisting}

%------------------------------------------------------------------------------
\subsection{Controller (GRASP)}

\begin{quote}
\textit{``Un Controller è il primo oggetto oltre il layer UI che riceve e coordina le operazioni di sistema.''}
\end{quote}

\subsubsection{Valutazione per Classe}

\begin{table}[h]
\centering
\begin{tabular}{|l|c|p{8cm}|}
\hline
\textbf{Classe} & \textbf{Voto} & \textbf{Motivazione} \\
\hline
CreaOrdineController & \ok & Coordina senza implementare business logic direttamente \\
\hline
CreaOrdineFacade & \ok & Pattern Facade per semplificare accesso (non un Controller vero) \\
\hline
CreaOrdineGUIController & \ok & Controller grafico (Boundary), gestisce solo UI \\
\hline
UsaVoucherController & \ok & Controller specializzato per sotto-use-case \\
\hline
\end{tabular}
\caption{Controller - Valutazione}
\end{table}

\textbf{Esempio positivo} - CreaOrdineController delega invece di implementare:

\begin{lstlisting}[style=javastyle, caption={Controller che Delega}]
public VoucherBean applicaVoucher(String codiceVoucher) {
    // DELEGA a UsaVoucherController (non implementa direttamente)
    return voucherController.applicaVoucherAOrdine(ordineCorrente, codiceVoucher);
}

public List<FoodBean> getProdottiBaseDisponibili() {
    List<Food> foodBase = FoodLazyFactory.getInstance().getAllFoodBase();
    // DELEGA conversione a FoodMapper
    return FoodMapper.toBeanList(foodBase);
}
\end{lstlisting}

%------------------------------------------------------------------------------
\subsection{Low Coupling}

\begin{quote}
\textit{``Minimizza le dipendenze tra classi per ridurre l'impatto dei cambiamenti.''}
\end{quote}

\subsubsection{Matrice delle Dipendenze}

\begin{table}[h]
\centering
\footnotesize
\begin{tabular}{|l|c|c|c|c|c|c|}
\hline
& \rotatebox{90}{Ordine} & \rotatebox{90}{Food} & \rotatebox{90}{Voucher} & \rotatebox{90}{Bean} & \rotatebox{90}{Mapper} & \rotatebox{90}{DAO} \\
\hline
CreaOrdineGUI & - & - & - & \ok & - & - \\
CreaOrdineFacade & - & - & - & \ok & - & - \\
CreaOrdineController & \ok & \ok & - & \ok & \ok & - \\
UsaVoucherController & \ok & - & \ok & \ok & \ok & - \\
FoodMapper & - & \ok & - & \ok & - & - \\
\hline
\end{tabular}
\caption{Matrice dipendenze (meno = meglio)}
\end{table}

\subsubsection{Valutazione}

\begin{table}[h]
\centering
\begin{tabular}{|l|c|p{8cm}|}
\hline
\textbf{Classe} & \textbf{Voto} & \textbf{Motivazione} \\
\hline
CreaOrdineGUIController & \ok & Dipende solo da Facade e Bean (nessuna Entity!) \\
\hline
CreaOrdineFacade & \ok & Isola GUI dal Controller, gestisce sessione \\
\hline
CreaOrdineController & \warn & Dipende da LazyFactory, Mapper, Entity (necessario) \\
\hline
UsaVoucherController & \ok & Dipendenze minime e stabili \\
\hline
\end{tabular}
\caption{Low Coupling - Valutazione}
\end{table}

\textbf{Punto di forza}: La GUI \textbf{non conosce le Entity} - comunica solo tramite Bean!

%------------------------------------------------------------------------------
\subsection{High Cohesion}

\begin{quote}
\textit{``Le responsabilità di una classe devono essere fortemente correlate e focalizzate.''}
\end{quote}

\subsubsection{Valutazione per Classe}

\begin{table}[h]
\centering
\begin{tabular}{|l|c|c|p{6cm}|}
\hline
\textbf{Classe} & \textbf{Voto} & \textbf{Responsabilità} & \textbf{Note} \\
\hline
CreaOrdineController & \ok & 1 & Solo orchestrazione creazione ordine \\
\hline
UsaVoucherController & \ok & 1 & Solo gestione voucher \\
\hline
CreaOrdineGUIController & \warn & 2 & UI + alcuni helper metodi \\
\hline
FoodMapper & \ok & 1 & Solo conversione Food$\to$Bean \\
\hline
Ordine & \ok & 1 & Solo rappresentazione ordine \\
\hline
FoodFactory & \ok & 1 & Solo creazione Food \\
\hline
\end{tabular}
\caption{High Cohesion - Valutazione}
\end{table}

\textbf{Miglioramento effettuato}: Rimozione metodi di conversione dai Controller:

\begin{lstlisting}[style=javastyle, caption={Prima: Bassa Coesione}]
// CreaOrdineController aveva 2 responsabilita:
public class CreaOrdineController {
    // 1. Business Logic
    public boolean aggiungiProdottoAOrdine(...) { ... }
    
    // 2. Conversione (RIMOSSO!)
    private FoodBean convertFoodToBean(Food food) { ... }
}
\end{lstlisting}

%------------------------------------------------------------------------------
\subsection{Polymorphism}

\begin{quote}
\textit{``Usa il polimorfismo per gestire alternative basate sul tipo.''}
\end{quote}

\subsubsection{Applicazioni nel Sistema}

\begin{enumerate}
    \item \textbf{Pattern Decorator per Food}:
\begin{lstlisting}[style=javastyle]
public abstract class Food {
    public abstract double getCosto();
    public abstract int getDurata();
}

public class Cipolla extends FoodDecorator {
    @Override
    public double getCosto() {
        return food.getCosto() + 0.50; // Polimorfismo!
    }
}
\end{lstlisting}

    \item \textbf{Pattern Strategy per Voucher}:
\begin{lstlisting}[style=javastyle]
public interface Voucher {
    double calcolaSconto(double subtotale);
}

public class VoucherPercentuale implements Voucher {
    @Override
    public double calcolaSconto(double subtotale) {
        return subtotale * (percentuale / 100.0);
    }
}

public class VoucherFisso implements Voucher {
    @Override
    public double calcolaSconto(double subtotale) {
        return (subtotale >= minimoOrdine) ? importoSconto : 0;
    }
}
\end{lstlisting}
\end{enumerate}

\subsubsection{Valutazione}
\begin{table}[h]
\centering
\begin{tabular}{|l|c|p{8cm}|}
\hline
\textbf{Area} & \textbf{Voto} & \textbf{Motivazione} \\
\hline
Food + Decorator & \ok & Polimorfismo per aggiungere add-on dinamicamente \\
\hline
Voucher + Strategy & \ok & Polimorfismo per diversi tipi di sconto \\
\hline
StatoOrdine & \warn & Enum con switch - potrebbe usare State Pattern \\
\hline
\end{tabular}
\caption{Polymorphism - Valutazione}
\end{table}

%------------------------------------------------------------------------------
\subsection{Pure Fabrication}

\begin{quote}
\textit{``Classe inventata per ottenere alta coesione e basso accoppiamento, non parte del dominio.''}
\end{quote}

\subsubsection{Classi Pure Fabrication nel Sistema}

\begin{table}[h]
\centering
\begin{tabular}{|l|p{6cm}|c|}
\hline
\textbf{Classe} & \textbf{Giustificazione} & \textbf{Voto} \\
\hline
FoodMapper & Centralizza conversione, non è un concetto di dominio & \ok \\
\hline
VoucherMapper & Come sopra & \ok \\
\hline
OrdineMapper & Come sopra & \ok \\
\hline
FoodFactory & Incapsula creazione complessa con Decorator & \ok \\
\hline
SessionManager & Gestione sessioni (concetto tecnico) & \ok \\
\hline
AlertUtils & Utility UI, non dominio & \ok \\
\hline
\end{tabular}
\caption{Pure Fabrication - Classi identificate}
\end{table}

%------------------------------------------------------------------------------
\subsection{Indirection}

\begin{quote}
\textit{``Introduce un intermediario per ridurre l'accoppiamento diretto.''}
\end{quote}

\subsubsection{Applicazioni}

\begin{enumerate}
    \item \textbf{CreaOrdineFacade}: Intermediario tra GUI e Controller
    \item \textbf{Mapper}: Intermediario tra Entity e Bean
    \item \textbf{LazyFactory}: Intermediario tra Controller e DAO
\end{enumerate}

\begin{table}[h]
\centering
\begin{tabular}{|l|l|l|c|}
\hline
\textbf{Indirection} & \textbf{Da} & \textbf{A} & \textbf{Voto} \\
\hline
Facade & GUI & Controller & \ok \\
\hline
Mapper & Controller & Bean/Entity & \ok \\
\hline
LazyFactory & Controller & DAO & \ok \\
\hline
\end{tabular}
\caption{Indirection - Pattern identificati}
\end{table}

%------------------------------------------------------------------------------
\subsection{Protected Variations}

\begin{quote}
\textit{``Proteggi gli elementi da variazioni in altri elementi tramite interfacce stabili.''}
\end{quote}

\subsubsection{Applicazioni}

\begin{table}[h]
\centering
\begin{tabular}{|l|p{5cm}|p{5cm}|}
\hline
\textbf{Variazione} & \textbf{Protezione} & \textbf{Implementazione} \\
\hline
Tipo di Voucher & Interfaccia \texttt{Voucher} & Strategy Pattern \\
\hline
Add-on Food & Classe astratta \texttt{Food} & Decorator Pattern \\
\hline
Tipo Persistenza & \texttt{DAOInterface} & Abstract Factory \\
\hline
Tipo UI & \texttt{Facade} & Facade Pattern \\
\hline
\end{tabular}
\caption{Protected Variations - Implementazioni}
\end{table}

%==============================================================================
\section{Riepilogo Valutazione GRASP}

\begin{table}[h]
\centering
\begin{tabular}{|l|c|c|c|c|c|c|c|c|c|}
\hline
\textbf{Classe} & \rotatebox{90}{Info Expert} & \rotatebox{90}{Creator} & \rotatebox{90}{Controller} & \rotatebox{90}{Low Coupling} & \rotatebox{90}{High Cohesion} & \rotatebox{90}{Polymorphism} & \rotatebox{90}{Pure Fab.} & \rotatebox{90}{Indirection} & \rotatebox{90}{Protected Var.} \\
\hline
CreaOrdineGUI & \ok & - & \ok & \ok & \warn & - & - & - & \ok \\
\hline
CreaOrdineFacade & - & \ok & - & \ok & \ok & - & - & \ok & \ok \\
\hline
CreaOrdineCtrl & \ok & \warn & \ok & \warn & \ok & - & - & - & - \\
\hline
UsaVoucherCtrl & \ok & - & \ok & \ok & \ok & - & - & - & - \\
\hline
Ordine & \ok & - & - & \ok & \ok & - & - & - & - \\
\hline
Food & \ok & - & - & \ok & \ok & \ok & - & - & \ok \\
\hline
Voucher & \ok & - & - & \ok & \ok & \ok & - & - & \ok \\
\hline
FoodMapper & - & - & - & \ok & \ok & - & \ok & \ok & - \\
\hline
FoodFactory & - & \ok & - & \ok & \ok & - & \ok & - & \ok \\
\hline
\end{tabular}
\caption{Matrice GRASP completa (\ok = conforme, \warn = parziale, - = non applicabile)}
\end{table}

%==============================================================================
\section{Problemi Identificati e Raccomandazioni}

\subsection{Problemi Minori}

\begin{enumerate}
    \item \textbf{CreaOrdineGUIController} (485 LOC):
    \begin{itemize}
        \item File grande con molte responsabilità UI
        \item \textbf{Raccomandazione}: Valutare estrazione di helper class per setup UI
    \end{itemize}
    
    \item \textbf{Magic Strings in FoodFactory}:
    \begin{itemize}
        \item Switch su stringhe letterali per tipi Food
        \item \textbf{Raccomandazione}: Usare enum o registry pattern
    \end{itemize}
    
    \item \textbf{StatoOrdine come Enum}:
    \begin{itemize}
        \item Gestito con switch invece di State Pattern
        \item \textbf{Raccomandazione}: Se la logica di stato cresce, valutare refactoring
    \end{itemize}
\end{enumerate}

\subsection{Punti di Forza}

\begin{enumerate}
    \item \textbf{Separazione BCE ben implementata}
    \item \textbf{Pattern Decorator} per Food con add-on
    \item \textbf{Pattern Strategy} per Voucher
    \item \textbf{Mapper Pattern} per disaccoppiamento Entity/Bean
    \item \textbf{Facade Pattern} per semplificare accesso
    \item \textbf{Delega voucher} a controller specializzato
\end{enumerate}

%==============================================================================
\section{Conclusioni}

L'analisi GRASP del caso d'uso ``Crea Ordine'' evidenzia un design \textbf{complessivamente maturo} con:

\begin{itemize}
    \item \textbf{9/9 principi GRASP} applicati correttamente o parzialmente
    \item \textbf{6 classi Pure Fabrication} che migliorano coesione e accoppiamento
    \item \textbf{2 pattern GoF} (Decorator, Strategy) ben integrati
    \item \textbf{Separazione chiara} tra layer Boundary, Control, Entity
\end{itemize}

\vspace{0.5cm}

\noindent\textbf{Valutazione complessiva}: \textcolor{okgreen}{\textbf{8.5/10}}

\vspace{0.3cm}

\noindent Il sistema rispetta i principi GRASP con solo lievi aree di miglioramento che non impattano significativamente la manutenibilità.

\end{document}
